\documentclass[11pt,a4paper]{article}
\usepackage[margin=2cm]{geometry}
\usepackage{enumitem}
\usepackage{booktabs}
\usepackage{xcolor}
\usepackage{titlesec}
\usepackage{hyperref}
\usepackage{amsmath,amssymb}
\usepackage{tcolorbox}
\usepackage{fancyhdr}

\definecolor{writeblue}{RGB}{30,80,160}
\definecolor{warnred}{RGB}{180,40,40}
\definecolor{tipgreen}{RGB}{30,120,50}
\definecolor{lightblue}{RGB}{230,240,255}
\definecolor{lightyellow}{RGB}{255,250,230}
\definecolor{lightred}{RGB}{255,235,235}

\tcbuselibrary{skins,breakable}

\newtcolorbox{writebox}[1]{
  colback=lightblue, colframe=writeblue, fonttitle=\bfseries,
  title={HANDWRITE: #1}, breakable
}

\newtcolorbox{whybox}[1]{
  colback=lightyellow, colframe=tipgreen!80!black, fonttitle=\bfseries,
  title={WHY THIS MATTERS: #1}, breakable
}

\newtcolorbox{dangerbox}[1]{
  colback=lightred, colframe=warnred, fonttitle=\bfseries,
  title={DANGER ZONE: #1}, breakable
}

\titleformat{\section}{\Large\bfseries\color{writeblue}}{Lecture \thesection:}{0.5em}{}[\titlerule]
\titleformat{\subsection}{\large\bfseries}{}{0em}{}

\pagestyle{fancy}
\fancyhf{}
\rhead{AC --- Handwriting Guide}
\lhead{Jawad M.K. Qayyum}
\rfoot{\thepage}
\setlength{\headheight}{14pt}

\begin{document}

\begin{center}
{\Huge\bfseries AC Handwriting Guide}\\[0.5em]
{\large What to write by hand, lecture by lecture}\\[0.3em]
{\large\color{warnred} The hand is the memory. LaTeX is the map. Paper is the learning.}
\end{center}

\vspace{0.5em}

\begin{tcolorbox}[colback=lightyellow, colframe=black!50]
\textbf{How to use this guide:}
\begin{enumerate}[nosep]
\item Read the lecture slides first
\item Open this guide to the matching lecture
\item \textbf{Handwrite} everything in the blue boxes on paper --- no typing, no copying
\item After handwriting, attempt the exercise \textbf{without} looking at the solution
\item Mark what you got wrong in your \texttt{error-catalog/}
\item Whatever you kept needing to look up $\rightarrow$ candidate for your A4 exam page
\end{enumerate}
\end{tcolorbox}

%=============================================================================
\section{Introduction and Turing Machines}
%=============================================================================

\begin{whybox}{Foundation of everything}
This lecture defines \textbf{what computation is}. Every later lecture builds on Turing machines. If you don't internalize the TM definition and how it processes input, nothing else will make sense.
\end{whybox}

\begin{writebox}{Turing Machine definition}
Handwrite the full formal definition of a Turing Machine:
\begin{itemize}[nosep]
\item $M = (Q, \Sigma, \Gamma, \delta, q_0, q_{accept}, q_{reject})$
\item What each component means (in your own words)
\item The transition function: $\delta: Q \times \Gamma \rightarrow Q \times \Gamma \times \{L, R\}$
\item Draw the tape, head, and state diagram for a simple example
\end{itemize}
\end{writebox}

\begin{writebox}{Key concepts}
Handwrite definitions and examples for:
\begin{itemize}[nosep]
\item \textbf{Configuration}: $(q, tape, head\ position)$ --- draw one
\item \textbf{Accepting vs.\ rejecting vs.\ looping}: what happens in each case
\item \textbf{Language of a TM}: $L(M) = \{w \mid M \text{ accepts } w\}$
\item \textbf{Computable language}: there exists a \textbf{halting} DTM that decides it
\item \textbf{Computably-enumerable (c.e.)}: there exists a TM that accepts it (may loop on non-members)
\item The difference: computable $\subset$ c.e.\ --- draw a Venn diagram
\end{itemize}
\end{writebox}

\begin{writebox}{Runtime analysis example}
Handwrite the primality-testing algorithm from the slides:
\begin{itemize}[nosep]
\item Write the code: \texttt{mystery(n)} with the $\sqrt{n}$ loop
\item Analyze its runtime step by step: why $\mathcal{O}(\sqrt{n})$?
\item Practice: what if the loop went to $n$ instead of $\sqrt{n}$? What changes?
\end{itemize}
\end{writebox}

\begin{dangerbox}{Your gap from Semester 1}
You scored 4 in Theoretical Foundations. That course covered basic TMs. If the definition above feels unfamiliar, spend extra time here. Draw 3 different TMs by hand --- one that accepts $\{0^n1^n\}$, one that copies a string, one that checks palindromes.
\end{dangerbox}

\subsection*{Exercise 1: What it tests}
After handwriting the above, attempt Exercise 1. It likely tests: constructing TMs for specific languages, tracing TM execution, and distinguishing computable vs.\ c.e.

%=============================================================================
\section{The Church-Turing Thesis}
%=============================================================================

\begin{whybox}{Why this matters}
This lecture says: all reasonable models of computation are equivalent. Multi-tape TMs, nondeterministic TMs --- they all recognize the same languages. This is why we can use TMs as \textit{the} model.
\end{whybox}

\begin{writebox}{Multi-tape Turing Machines}
Handwrite:
\begin{itemize}[nosep]
\item How a multi-tape TM differs from a single-tape TM (draw it)
\item The simulation theorem: a $k$-tape TM can be simulated by a single-tape TM
\item The cost: if the $k$-tape TM runs in $t$ steps, the single-tape simulation takes $\mathcal{O}(t^2)$ steps
\item Why: the single tape must sweep back and forth to simulate multiple tapes
\end{itemize}
\end{writebox}

\begin{writebox}{Nondeterministic Turing Machines (NTMs)}
Handwrite:
\begin{itemize}[nosep]
\item Definition: $\delta: Q \times \Gamma \rightarrow \mathcal{P}(Q \times \Gamma \times \{L,R\})$ (set of possible transitions)
\item NTM accepts if \textbf{any} computation branch accepts (draw the tree)
\item Equivalence: NTMs and DTMs recognize the same class of languages
\item The simulation cost: a NTM running in $t$ steps can be simulated by a DTM in $2^{\mathcal{O}(t)}$ steps
\item Draw the computation tree: branching paths, one accepting $=$ accept
\end{itemize}
\end{writebox}

\begin{writebox}{The Church-Turing Thesis}
Handwrite in your own words:
\begin{itemize}[nosep]
\item Statement: ``Every effectively computable function is computable by a Turing machine''
\item This is a \textbf{thesis}, not a theorem --- it cannot be proven, only supported by evidence
\item Evidence: every proposed model of computation (lambda calculus, RAM machines, etc.) has been shown equivalent to TMs
\end{itemize}
\end{writebox}

\subsection*{Exercise 2: What it tests}
Likely: simulating multi-tape TMs, tracing NTM computation trees, proving equivalences between models.

%=============================================================================
\section{Computability}
%=============================================================================

\begin{whybox}{The hierarchy of languages}
This lecture maps out what can and cannot be computed. The Venn diagram (regular $\subset$ context-free $\subset$ computable $\subset$ c.e.) is fundamental. You need to know examples in each region.
\end{whybox}

\begin{writebox}{The language hierarchy}
Handwrite the full Venn diagram with examples:
\begin{itemize}[nosep]
\item \textbf{Regular}: $\{w \mid w \text{ contains } 01\}$ --- recognized by DFA
\item \textbf{Context-free}: $\{w\#w^R\}$ --- recognized by PDA
\item \textbf{Computable (decidable)}: $\{w\#w\}$ --- requires TM, but always halts
\item \textbf{Computably-enumerable}: $A_{TM} = \{\langle M,w \rangle \mid M \text{ accepts } w\}$
\item \textbf{Not c.e.}: $\overline{A_{TM}}$
\end{itemize}
\end{writebox}

\begin{writebox}{The Halting Problem and undecidability}
Handwrite the full proof that $A_{TM}$ is undecidable:
\begin{itemize}[nosep]
\item Assume $H$ decides $A_{TM}$
\item Construct $D$: on input $\langle M \rangle$, run $H(\langle M, \langle M \rangle \rangle)$, then do the opposite
\item Run $D$ on $\langle D \rangle$ --- contradiction (accepts iff rejects)
\item Therefore no such $H$ exists
\item \textbf{Write this proof 3 times by hand until you can do it from memory}
\end{itemize}
\end{writebox}

\begin{writebox}{Closure properties}
Handwrite:
\begin{itemize}[nosep]
\item Computable languages are closed under: union, intersection, complement, concatenation, Kleene star
\item C.e.\ languages are closed under: union, intersection, concatenation, Kleene star
\item C.e.\ languages are \textbf{NOT} closed under: complement
\item Key theorem: $L$ is computable $\iff$ both $L$ and $\overline{L}$ are c.e.
\end{itemize}
\end{writebox}

\begin{dangerbox}{The diagonalization proof}
The proof of undecidability of $A_{TM}$ uses diagonalization. This is the hardest concept to internalize. If you cannot reproduce this proof from memory, you are not ready for the exam. Handwrite it repeatedly.
\end{dangerbox}

\subsection*{Exercise 3: What it tests}
Proving languages are (un)decidable, closure property arguments, constructing TMs for decidable languages.

%=============================================================================
\section{Reductions}
%=============================================================================

\begin{whybox}{The core exam skill}
Reductions are how you prove new problems are undecidable by relating them to known undecidable problems. This is tested in almost every exam. Master this.
\end{whybox}

\begin{writebox}{What is a reduction?}
Handwrite:
\begin{itemize}[nosep]
\item $A \leq_m B$ means: there exists a computable function $f$ such that $w \in A \iff f(w) \in B$
\item Intuition: ``$A$ is no harder than $B$'' --- if you can solve $B$, you can solve $A$
\item Contrapositive: if $A$ is undecidable and $A \leq_m B$, then $B$ is undecidable
\item Draw the diagram: input $w \xrightarrow{f} f(w)$, then decide $f(w) \in B$
\end{itemize}
\end{writebox}

\begin{writebox}{Key reductions to memorize}
Handwrite the full reduction for each:
\begin{itemize}[nosep]
\item $A_{TM} \leq_m HALT_{TM}$: prove the halting problem is undecidable
\item $A_{TM} \leq_m E_{TM}$: prove emptiness testing is undecidable ($E_{TM} = \{\langle M \rangle \mid L(M) = \emptyset\}$)
\item $A_{TM} \leq_m EQ_{TM}$: prove equivalence testing is undecidable
\item For each: write the mapping function $f$, explain why $w \in A_{TM} \iff f(w) \in B$
\end{itemize}
\end{writebox}

\begin{writebox}{Rice's Theorem}
Handwrite:
\begin{itemize}[nosep]
\item Statement: every non-trivial property of c.e.\ languages is undecidable
\item ``Non-trivial'' means: some TMs have the property and some don't
\item Examples of properties: ``$L(M) = \emptyset$'', ``$L(M)$ is finite'', ``$L(M)$ is regular''
\item All undecidable by Rice's theorem
\item Proof sketch: reduction from $A_{TM}$
\end{itemize}
\end{writebox}

\subsection*{Exercise 4: What it tests}
Constructing reductions between problems, applying Rice's theorem, determining if a language property is decidable.

%=============================================================================
\section{Complexity Theory (P and NP)}
%=============================================================================

\begin{whybox}{From ``can we compute it?'' to ``can we compute it fast?''}
Lectures 1--4 asked what is computable. Now we ask: among computable problems, which ones can be solved \textbf{efficiently}? This is where P, NP, and NP-completeness live.
\end{whybox}

\begin{writebox}{Complexity classes}
Handwrite definitions:
\begin{itemize}[nosep]
\item $\text{TIME}(t(n))$: languages decidable by a DTM in $\mathcal{O}(t(n))$ steps
\item $\textbf{P} = \bigcup_{k \geq 1} \text{TIME}(n^k)$: polynomial time on a DTM
\item $\text{NTIME}(t(n))$: languages decidable by a NTM in $\mathcal{O}(t(n))$ steps
\item $\textbf{NP} = \bigcup_{k \geq 1} \text{NTIME}(n^k)$: polynomial time on a NTM
\item $\textbf{P} \subseteq \textbf{NP}$ (every DTM is trivially a NTM)
\item The big open question: $\textbf{P} \stackrel{?}{=} \textbf{NP}$
\end{itemize}
\end{writebox}

\begin{writebox}{NP via verifiers}
Handwrite the equivalent definition:
\begin{itemize}[nosep]
\item $L \in \textbf{NP}$ $\iff$ there exists a polynomial-time verifier $V$ such that:
\item $w \in L \iff \exists$ certificate $c$ with $|c| \leq p(|w|)$ and $V(w,c)$ accepts
\item Example: SAT --- certificate is an assignment, verifier checks it in polynomial time
\item Example: CLIQUE --- certificate is a set of $k$ vertices, verifier checks edges
\end{itemize}
\end{writebox}

\begin{writebox}{Examples in P and NP}
Handwrite 3 problems in P and 3 in NP with brief justification:
\begin{itemize}[nosep]
\item \textbf{In P}: PATH (graph reachability), 2-SAT, sorting
\item \textbf{In NP (not known in P)}: SAT, CLIQUE, HAMPATH, VERTEX-COVER, SUBSET-SUM
\end{itemize}
\end{writebox}

\subsection*{Exercise 5: What it tests}
Classifying problems into P/NP, constructing verifiers, understanding certificates.

%=============================================================================
\section{NP-Completeness}
%=============================================================================

\begin{dangerbox}{The most exam-critical lecture}
NP-completeness and polynomial reductions are the heart of AC exams. You must be able to both \textbf{prove} a problem is NP-complete and \textbf{construct} a reduction.
\end{dangerbox}

\begin{writebox}{Polynomial-time reduction}
Handwrite:
\begin{itemize}[nosep]
\item $A \leq_p B$: there exists a \textbf{polynomial-time} computable function $f$ such that $w \in A \iff f(w) \in B$
\item If $A \leq_p B$ and $B \in \textbf{P}$, then $A \in \textbf{P}$
\item If $A \leq_p B$ and $A$ is NP-hard, then $B$ is NP-hard
\end{itemize}
\end{writebox}

\begin{writebox}{NP-complete definition}
Handwrite:
\begin{itemize}[nosep]
\item $L$ is \textbf{NP-hard} if $\forall A \in \textbf{NP}: A \leq_p L$
\item $L$ is \textbf{NP-complete} if $L \in \textbf{NP}$ AND $L$ is NP-hard
\item To prove NP-completeness: (1) show $L \in$ NP, (2) reduce a known NP-complete problem to $L$
\end{itemize}
\end{writebox}

\begin{writebox}{The Cook-Levin Theorem}
Handwrite:
\begin{itemize}[nosep]
\item Statement: SAT is NP-complete
\item This is the ``first'' NP-complete problem --- proved directly, not by reduction
\item Once SAT is NP-complete, we can prove others by reducing FROM SAT
\end{itemize}
\end{writebox}

\begin{writebox}{The reduction chain --- MEMORIZE THIS}
Handwrite the full chain:
\begin{center}
SAT $\leq_p$ 3-SAT $\leq_p$ CLIQUE $\leq_p$ VERTEX-COVER $\leq_p$ \ldots
\end{center}
For each reduction, handwrite:
\begin{itemize}[nosep]
\item The mapping: how do you transform an instance of $A$ into an instance of $B$?
\item The proof: why does $w \in A \iff f(w) \in B$?
\item Draw a small example for each
\end{itemize}
\end{writebox}

\subsection*{Exercise 6: What it tests}
Proving NP-completeness via reduction, constructing polynomial-time mappings, showing both directions of the iff.

%=============================================================================
\section{More NP-Completeness and Beyond}
%=============================================================================

\begin{writebox}{Additional NP-complete problems}
Handwrite a reduction proof for each:
\begin{itemize}[nosep]
\item SUBSET-SUM is NP-complete
\item HAMPATH (Hamiltonian Path) is NP-complete
\item 3-COLORING is NP-complete
\item For each: what known NP-complete problem do you reduce FROM? What is the mapping?
\end{itemize}
\end{writebox}

\begin{writebox}{coNP and the complexity landscape}
Handwrite:
\begin{itemize}[nosep]
\item $\textbf{coNP} = \{\overline{L} \mid L \in \textbf{NP}\}$
\item $\textbf{P} \subseteq \textbf{NP} \cap \textbf{coNP}$
\item Draw the Venn diagram: P inside NP, P inside coNP, NP-complete at the boundary
\item If any NP-complete problem is in P, then $\textbf{P} = \textbf{NP}$
\end{itemize}
\end{writebox}

\begin{writebox}{Space complexity (if covered)}
Handwrite:
\begin{itemize}[nosep]
\item PSPACE: problems solvable with polynomial space
\item $\textbf{P} \subseteq \textbf{NP} \subseteq \textbf{PSPACE} \subseteq \textbf{EXPTIME}$
\item PSPACE-complete examples: TQBF, generalized games
\item Savitch's theorem: $\text{NSPACE}(f(n)) \subseteq \text{DSPACE}(f(n)^2)$
\end{itemize}
\end{writebox}

\subsection*{Exercise 7: What it tests}
More complex reductions, possibly space complexity, understanding the full complexity landscape.

%=============================================================================
\vspace{1em}
\section*{Your A4 Page --- What Goes On It}
%=============================================================================

\begin{tcolorbox}[colback=lightred, colframe=warnred, fonttitle=\bfseries, title={After handwriting all 7 lectures, these are A4 candidates}]
Only put on the A4 what you \textbf{cannot hold in memory} after handwriting everything above:
\begin{enumerate}[nosep]
\item TM formal definition ($M = (Q, \Sigma, \Gamma, \delta, q_0, q_{acc}, q_{rej})$)
\item The reduction chain: SAT $\leq_p$ 3-SAT $\leq_p$ CLIQUE $\leq_p$ VERTEX-COVER $\leq_p$ \ldots
\item Complexity class inclusions: P $\subseteq$ NP $\subseteq$ PSPACE $\subseteq$ EXPTIME
\item Rice's Theorem statement
\item Closure properties table (computable vs.\ c.e.)
\item Key reduction mappings you keep forgetting
\item NTM simulation costs ($2^{\mathcal{O}(t)}$ for DTM, $\mathcal{O}(t^2)$ for multi-tape)
\item The diagonalization proof skeleton (if you still can't reproduce it)
\end{enumerate}
If you can reproduce all of the above from memory after handwriting --- your A4 page is free for edge cases.
\end{tcolorbox}

%=============================================================================
\section*{Error Catalog Template}
%=============================================================================

\begin{tcolorbox}[colback=lightyellow, colframe=black!50]
After each exercise, add entries to \texttt{AC/error-catalog/} in this format:
\begin{verbatim}
Exercise: 3, Question: 2b
What I got wrong: forgot that c.e. is NOT closed under complement
Why: confused computable closure with c.e. closure
Fix: computable = closed under everything, c.e. = NOT complement
A4 candidate: YES - add closure properties table
\end{verbatim}
\end{tcolorbox}

\end{document}
